\definecolor{myBlue}{RGB}{190,203,211}
\definecolor{myRed}{RGB}{244,186,175}
\definecolor{myGreen}{RGB}{114,138,122}
\definecolor{myYellow}{RGB}{241,208,157}

\section{The Multi-level Weighted Essentially Non-Oscillatory (WENO) Method}
We develop a accurate, versatile, robust and efficient finite volume 
weighted essentially non oscillatory (WENO) techniques for solving energy and composition 
conservation laws \cite{MultiLevelWeno}. 

\subsection{Stencil Polynomial}
Given a stencil $S = \{E_0, E_2, \dots, E_k\}$ of the quasiuniform logically rectangular mesh. We can define
an associated tensor product stencil polynomial 
$Q_{r,s} \in \mathbb{Q}_{(r,s)} \vcentcolon= \mathbb{P}_{r-1}(x) \otimes \mathbb{P}_{s-1}(y)$,
where $r$ and $s$ equal to the stencil sizes in the 
x- and y-directions, respectively.
We scale our stencil polynomial individually as
\begin{equation} \label{eq:stencilpolynomial}
		  Q_{r,s} = \sum_{p<r, q<s} c_{p,q} \Big(\frac{x - x_S}{h_S}\Big)^p 
		  \Big(\frac{y - y_S}{h_S}\Big)^q,
\end{equation}
where $(x_S,y_S)$ is the center point of the stencil 
$\Omega_S = \cup_{E\in S} \subset \mathbb{R}^2$, and $h_S = 
\sqrt{\sum_{k}E_k/k}$ is the diameter of 
average cell size.
In finite volume scenerio, we are dealing with cell averaged value of the
scalar transport variable
\begin{equation}\label{eq:cellaveraged}
		  \overline{u}_E = \frac{1}{|E|}\int_E u(\mathbf{x}) d\mathbf{x}.
\end{equation}
We can calculate coefficents $c_{p,q}$ by imposing a constraint that cell averaged
value \eqref{eq:cellaveraged} is preserved through integral with stencil polynomial
\begin{equation}\label{eq:stencilpolycoefficent}
		  \frac{1}{|E_k|} \int_{E_k} Q(\mathbf{x}) d\mathbf{x} = \overline{u}_{E_k}, \quad 
		  \forall E_k \in S.
\end{equation}
I will replace $\overline{u}_{E_k}$ with $\overline{u}_k$ in the following
discussion for simplicity.

We can represent stencil polynomial with basis polynomials 
$Q_{r,s} = \sum_{k} \overline{u}_{k} \phi^k_{r,s}$,
where basis polynomials can be calcualted as
\begin{equation}\label{eq:basispoly}
		  \frac{1}{|E_{k'}|} \int_{E_{k'}} \phi^k_{r,s} (\mathbf{x}) d\mathbf{x} = \left \{
					 \begin{aligned}
								&1 \quad  k' = k\\
								&0 \quad  k' \neq k
					 \end{aligned}
					 \right. \quad \forall (E_{k'}, E_k) \in S,
\end{equation}
which would lead to a $(r \times s) \times (r \times s)$ non-singular linear system.


\subsection{ML-WENO weighting}
Given a target cell $E_{i,j} \in \mathbb{\tau}_h$, we can create a nonempty set of stencils $S_l \ni E_{i,j}, l = 0,1, \dots, L$.
Each stencil will 
\begin{figure}[ht]
\centering
\input chapters/mlstencil.tex
\caption{An example of (3,2,1) multi-level stencils combination on logically rectangular 
quadrilateral mesh.\label{fig:quadmesh}}
\end{figure}


\subsection{Boundary Treatment}

\subsection{Time Stepping}
